\documentclass[11pt]{article}
\usepackage[utf8]{inputenc}
\usepackage[margin=1in]{geometry}
\usepackage{amsmath}
\usepackage{graphicx}
\usepackage{booktabs}
\usepackage{hyperref}
\usepackage[round]{natbib}
\usepackage{float}
\usepackage{multirow}
\usepackage{tabularx}

\title{A SUMO E3 Ligase Promotes Long Non-Coding RNA Transcription to Regulate Small RNA-Directed DNA Elimination}
\author{}
\date{}

\begin{document}
\maketitle

\begin{abstract}
Small RNA-directed chromatin silencing requires nascent long non-coding RNA (lncRNA) transcription from target loci, but the molecular mechanisms promoting this transcription remain poorly understood. Using the ciliate \textit{Tetrahymena thermophila}, where the source (germline micronucleus, MIC) and target (somatic macronucleus, MAC) of small RNAs reside in physically separate nuclei, we identified Ema2 (TTHERM\_00113330) as a conjugation-specific SUMO E3 ligase that promotes genome-wide lncRNA transcription from the parental MAC. Ema2 contains a conserved SP-RING domain, is exclusively expressed during conjugation (3--14 hours post-mixing, hpm), and localizes to the MAC with a parental-to-new MAC switch at $\sim$8 hpm. EMA2 knockout partially blocks DNA elimination of $\sim$12,000 internal eliminated sequences (IESs) and produces inviable sexual progeny. Mechanistically, Ema2 interacts directly with the SUMO E2 enzyme Ubc9, promotes SUMOylation of the transcription regulator Spt6 (identified by mass spectrometry of His-Smt3-conjugated proteins), and is required for the chromatin association of both Spt6 and RNA polymerase II (Rpb3 subunit). Loss of Ema2 abolishes lncRNA accumulation from the parental MAC (assessed by RT-PCR spanning MDS-IES junctions and J2 anti-dsRNA immunofluorescence), thereby blocking target-directed scan RNA degradation (TDSD). However, SUMOylation-defective Spt6 mutants do not phenocopy the lncRNA defect, indicating that Ema2 has additional critical SUMOylation targets beyond Spt6. These findings establish SUMOylation as a key mechanism connecting lncRNA transcription to small RNA-mediated genome surveillance.
\end{abstract}

\section{Introduction}

Small RNA-directed chromatin silencing is a conserved mechanism for genome defense across eukaryotes \citep{malone2009small, czech2018pirna, holoch2015rnai}. In diverse organisms---including fission yeast, \textit{Drosophila}, and plants---small RNAs guide effector complexes to target chromatin through base-pairing with nascent transcripts from the target loci. This creates a paradox: regions destined for transcriptional silencing must first be actively transcribed to produce the nascent RNA scaffold required for small RNA recognition \citep{holoch2015rnai}. In most model systems, the source and target loci of small RNAs overlap spatially and temporally, complicating efforts to dissect the mechanism of lncRNA transcription independently from the silencing outcome.

The ciliated protozoan \textit{Tetrahymena thermophila} provides a uniquely tractable system for studying this problem \citep{chalker2013dna, kataoka2010twi1}. \textit{Tetrahymena} maintains two functionally distinct nuclei: a germline micronucleus (MIC, $\sim$200 Mb) that is transcriptionally silent during vegetative growth and serves as the genetic reservoir, and a somatic macronucleus (MAC, $\sim$103 Mb) that is transcriptionally active and determines the cell's phenotype \citep{eisen2006macronuclear}. During sexual conjugation, the MIC undergoes meiosis to produce the new MAC for the progeny cell. This process requires programmed DNA elimination of approximately 12,000 internal eliminated sequences (IESs), removing $\sim$97 Mb of MIC-specific DNA \citep{yao2003dna, chalker2005epigenetics}.

DNA elimination is guided by scan RNAs (scnRNAs), $\sim$26--31 nucleotide small RNAs produced from the MIC genome during early conjugation and loaded onto the Piwi protein Twi1 \citep{mochizuki2002analysis, mochizuki2004conjugation}. In a process termed target-directed scnRNA degradation (TDSD), scnRNAs complementary to MAC-destined sequences (MDS) are selectively degraded upon scanning lncRNA transcripts from the parental MAC, while IES-matching scnRNAs are retained \citep{aronica2008studying, schoeberl2012twi1, noto2015small}. The retained IES-matching scnRNAs then guide H3K27me3 and H3K9me3 deposition on corresponding IES sequences in the developing new MAC \citep{liu2007rnapii}, triggering their excision.

A critical but poorly understood step in this pathway is the active transcription of lncRNAs from the parental MAC during mid-conjugation ($\sim$3--8 hpm). These lncRNAs serve as the scanning substrate for TDSD and must be produced genome-wide to ensure comprehensive surveillance. While several components of the downstream pathway have been characterized, the molecular mechanism that promotes this lncRNA transcription was unknown.

SUMOylation---the conjugation of Small Ubiquitin-like Modifier (SUMO) proteins to target substrates---regulates diverse nuclear processes including transcription, chromatin remodeling, and DNA repair \citep{johnson2005role, gareau2010sumo}. Recent work in \textit{Drosophila} identified the SUMO E3 ligase Su(var)2-10 as a link between piRNA-guided target recognition and chromatin silencing \citep{ninova2020suvar}. However, whether SUMOylation plays a role in the lncRNA transcription step of small RNA pathways had not been established.

Here, we identify Ema2 as a conjugation-specific SUMO E3 ligase in \textit{Tetrahymena} that promotes lncRNA transcription from the parental MAC, connecting the SUMOylation machinery to the scnRNA-mediated genome surveillance pathway.

\section{Results}

\subsection{Ema2 Is a Conjugation-Specific Protein That Localizes to the Macronucleus}

Analysis of microarray expression data \citep{miao2009microarray} revealed that EMA2 mRNA is exclusively expressed during conjugation, with transcript levels undetectable in vegetative or starved cells and peaking between 3 and 6 hpm (Table~\ref{tab:expression}).

\begin{table}[H]
\centering
\caption{EMA2 mRNA expression profile across growth and conjugation stages, based on microarray data from \citet{miao2009microarray}. Expression scored as undetectable (--), low (+), or high (++).}
\label{tab:expression}
\begin{tabular}{lc}
\toprule
\textbf{Stage / Condition} & \textbf{EMA2 mRNA Level} \\
\midrule
Growing (low density) & -- \\
Growing (medium density) & -- \\
Growing (high density) & -- \\
Starved (0--24 hr) & -- \\
Conjugation 0 hpm & + \\
Conjugation 3--6 hpm & ++ \\
Conjugation 8--12 hpm & + \\
Conjugation 14--18 hpm & $\pm$ \\
\bottomrule
\end{tabular}
\end{table}

Immunofluorescence of endogenously HA-tagged Ema2 showed that the protein first appears in the parental MAC at 3 hpm, persists through 6 hpm, then switches to the developing new MAC at $\sim$8 hpm before fading by 14 hpm (Table~\ref{tab:localization}). This parental-to-new MAC transition is reminiscent of the Piwi protein Twi1 and the PRC2 complex, placing Ema2 at the right time and place for a role in the scnRNA pathway.

\begin{table}[H]
\centering
\caption{Ema2-HA protein localization during conjugation, assessed by immunofluorescence with anti-HA antibody (1:1,000 dilution). pMAC = parental macronucleus; nMAC = new macronucleus.}
\label{tab:localization}
\begin{tabular}{lcl}
\toprule
\textbf{Time Point} & \textbf{Location} & \textbf{Signal Intensity} \\
\midrule
Vegetative & Not detected & -- \\
3 hpm & pMAC & Strong \\
6 hpm & pMAC & Strong \\
8 hpm & nMAC (pMAC gone) & Moderate \\
12 hpm & nMAC & Weak \\
14 hpm & nMAC & Very weak \\
\bottomrule
\end{tabular}
\end{table}

\subsection{EMA2 Knockout Partially Blocks DNA Elimination and Abolishes Progeny Viability}

Somatic knockout of all MAC copies of EMA2 was achieved and confirmed by genomic PCR. DNA elimination was assessed at 36 hpm using FISH with probes against the Tlr1 transposable element, which is normally eliminated from the new MAC. In wild-type crosses, Tlr1 signal was restricted to the MIC, confirming complete elimination. In EMA2 KO crosses, Tlr1 signal was detected in both the MIC and the new MAC, indicating partial failure of DNA elimination (Table~\ref{tab:dna_elimination}). The FISH signal intensity in EMA2 KO new MACs was consistently lower than in TWI1 KO (complete elimination block), suggesting that Ema2 loss causes a partial rather than complete defect.

\begin{table}[H]
\centering
\caption{DNA elimination assessed by Tlr1 FISH at 36 hpm. Signal presence (+) or absence (--) scored in the micronucleus (MIC) and new macronucleus (nMAC). Relative intensity in nMAC graded on a 0--3 scale.}
\label{tab:dna_elimination}
\begin{tabular}{lcccc}
\toprule
\textbf{Genotype} & \textbf{MIC Signal} & \textbf{nMAC Signal} & \textbf{nMAC Intensity (0--3)} & \textbf{Viable Progeny} \\
\midrule
Wild-type & + & -- & 0 & Yes \\
EMA2 KO & + & + & 1.5 & No \\
TWI1 KO & + & + & 3.0 & No \\
\bottomrule
\end{tabular}
\end{table}

\subsection{Ema2 Is Required for Target-Directed scnRNA Degradation}

Northern blot analysis with the Mi-9 probe (complementary to an MDS region) revealed that TDSD is blocked in EMA2 KO cells. In wild-type crosses, MDS-complementary scnRNAs were progressively degraded from 3 to 6 hpm. In EMA2 KO crosses, these scnRNAs persisted through 6 hpm, demonstrating a failure of TDSD (Table~\ref{tab:tdsd}). Total scnRNA biogenesis was normal in EMA2 KO, confirming that the defect is specific to the scanning/degradation step rather than scnRNA production.

\begin{table}[H]
\centering
\caption{Target-directed scnRNA degradation (TDSD) assessed by northern blot with Mi-9 (MDS-complementary) probe. Signal intensity scored on a 0--3 scale (0 = absent, 3 = maximum).}
\label{tab:tdsd}
\begin{tabular}{lccccl}
\toprule
\textbf{Genotype} & \textbf{3 hpm} & \textbf{4.5 hpm} & \textbf{6 hpm} & \textbf{scnRNA Biogenesis} & \textbf{TDSD Status} \\
\midrule
Wild-type & 3.0 & 1.5 & 0.5 & Normal & Functional \\
EMA2 KO & 3.0 & 2.8 & 2.5 & Normal & Blocked \\
EMA2 rescue (HA) & 3.0 & 1.8 & 0.8 & Normal & Restored \\
\bottomrule
\end{tabular}
\end{table}

H3K27me3 immunofluorescence in the developing new MAC was reduced or absent in EMA2 KO crosses compared to wild-type, consistent with impaired heterochromatin formation downstream of the TDSD defect.

\subsection{Ema2 Contains a Conserved SP-RING Domain and Functions as a SUMO E3 Ligase}

Sequence analysis revealed that Ema2 contains an SP-RING domain with conserved cysteine and histidine residues required for zinc coordination, homologous to known SUMO E3 ligases across eukaryotes (Table~\ref{tab:spring}).

\begin{table}[H]
\centering
\caption{SP-RING domain conservation across SUMO E3 ligases. Conserved Zn-binding residues (Cys and His) are present in all homologs.}
\label{tab:spring}
\begin{tabular}{llcc}
\toprule
\textbf{Organism} & \textbf{Protein} & \textbf{SP-RING Domain} & \textbf{Key Zn Residues} \\
\midrule
\textit{T. thermophila} & Ema2 & Present & Conserved \\
\textit{S. cerevisiae} & MMS21 & Present & Conserved \\
\textit{S. cerevisiae} & SIZ1 & Present & Conserved \\
\textit{D. melanogaster} & Su(var)2-10 & Present & Conserved \\
\textit{H. sapiens} & PIAS1 & Present & Conserved \\
\bottomrule
\end{tabular}
\end{table}

GST pulldown assays using recombinant GST-Ema2 and His-Ubc9 expressed in \textit{E.~coli} confirmed direct physical interaction between Ema2 and the SUMO E2 conjugating enzyme Ubc9. GST alone showed no interaction. In vitro SUMOylation assays demonstrated that addition of Ema2 to the E1 + E2 (Ubc9) reaction enhanced overall SUMOylation activity, confirming its E3 ligase function (Table~\ref{tab:interaction}).

\begin{table}[H]
\centering
\caption{In vitro interaction and SUMOylation assays for Ema2. GST pulldown assessed physical interaction with His-Ubc9; SUMOylation assays measured enhancement of E1+E2 activity.}
\label{tab:interaction}
\begin{tabular}{lcc}
\toprule
\textbf{Condition} & \textbf{Ubc9 Interaction} & \textbf{SUMOylation Activity} \\
\midrule
GST alone + His-Ubc9 & -- & N/A \\
GST-Ema2 + His-Ubc9 & + & N/A \\
E1 + E2 (Ubc9) only & N/A & 1.0$\times$ (basal) \\
E1 + E2 + Ema2 & N/A & 3.2$\times$ (enhanced) \\
\bottomrule
\end{tabular}
\end{table}

\subsection{Ema2 Promotes SUMOylation of the Transcription Regulator Spt6}

To identify Ema2-dependent SUMOylation targets \textit{in vivo}, His-tagged Smt3 (SUMO) was expressed in wild-type and EMA2 KO backgrounds, and His-Smt3-conjugated proteins were purified by Ni-NTA pulldown from conjugating cells at 6 hpm. Label-free quantification (LFQ) mass spectrometry identified Spt6 as an Ema2-dependent SUMOylation target, with higher LFQ intensity in wild-type relative to EMA2 KO. Western blot analysis of Spt6-HA during conjugation revealed slower-migrating species appearing at 4.5--6 hpm, consistent with SUMOylated forms (Table~\ref{tab:spt6_sumo}).

\begin{table}[H]
\centering
\caption{Spt6 SUMOylation during conjugation. Western blot of Spt6-HA showing presence (+) or absence (--) of slower-migrating (putatively SUMOylated) bands.}
\label{tab:spt6_sumo}
\begin{tabular}{lcc}
\toprule
\textbf{Time Point} & \textbf{Main Band} & \textbf{Slower-Migrating Band(s)} \\
\midrule
Vegetative & + & -- \\
0 hpm & + & -- \\
4.5 hpm & + & + \\
6 hpm & + & + \\
No-tag control & -- & -- \\
\bottomrule
\end{tabular}
\end{table}

\subsection{Ema2 Is Required for lncRNA Accumulation from the Parental MAC}

RT-PCR using primers spanning MDS-IES junctions---which detect lncRNAs reading through IES boundaries but not processed mRNAs---showed that lncRNA products were present in wild-type at 6 hpm but absent or strongly reduced in EMA2 KO (Table~\ref{tab:lncrna}). Genomic DNA controls confirmed that the target loci were intact in KO strains, ruling out genomic deletion as the cause of lncRNA loss.

\begin{table}[H]
\centering
\caption{lncRNA detection by RT-PCR and dsRNA immunofluorescence at 6 hpm. RT-PCR uses primers spanning MDS-IES junctions. J2 antibody detects long dsRNA in the parental MAC.}
\label{tab:lncrna}
\begin{tabular}{lccc}
\toprule
\textbf{Genotype} & \textbf{RT-PCR (MDS-IES)} & \textbf{J2 dsRNA Signal (pMAC)} & \textbf{Genomic DNA Intact?} \\
\midrule
Wild-type & + & Strong & Yes \\
EMA2 KO & -- & Absent/reduced & Yes \\
\bottomrule
\end{tabular}
\end{table}

\subsection{Ema2 Promotes Chromatin Association of Spt6 and RNA Polymerase II}

Immunofluorescence comparing wild-type and EMA2 KO crosses at mid-conjugation showed that both Spt6-HA and the RNAPII subunit Rpb3 had reduced chromatin association in EMA2 KO cells (Table~\ref{tab:chromatin}). Since Spt6 is a histone chaperone that facilitates transcription elongation \citep{yoh2007spt6, ardehali2009spt6}, and RNAPII is the enzyme responsible for lncRNA transcription, their loss from chromatin in EMA2 KO provides a mechanistic explanation for the lncRNA accumulation defect.

\begin{table}[H]
\centering
\caption{Chromatin association of Spt6 and RNAPII (Rpb3) during mid-conjugation, assessed by immunofluorescence. Signal scored as strong (+++), moderate (++), weak (+), or absent (--).}
\label{tab:chromatin}
\begin{tabular}{lccc}
\toprule
\textbf{Cross} & \textbf{Spt6-HA Signal} & \textbf{Rpb3 Signal} & \textbf{Time Point} \\
\midrule
WT $\times$ WT & +++ & +++ & Mid-conjugation \\
EMA2 KO $\times$ EMA2 KO & + & -- & Mid-conjugation \\
\bottomrule
\end{tabular}
\end{table}

\subsection{SUMOylation-Defective Spt6 Does Not Phenocopy EMA2 Knockout}

To test whether Spt6 SUMOylation alone accounts for Ema2's role in lncRNA transcription, we examined a SUMOylation-defective Spt6 mutant. Surprisingly, this mutant showed normal lncRNA production and normal DNA elimination, in contrast to the EMA2 KO phenotype (Table~\ref{tab:spt6_mutant}). This indicates that Ema2 has additional SUMOylation targets beyond Spt6 that are critical for lncRNA transcription, consistent with its E3 ligase activity potentially acting on multiple chromatin-associated substrates.

\begin{table}[H]
\centering
\caption{Comparison of EMA2 KO and SUMOylation-defective Spt6 mutant phenotypes for lncRNA production and DNA elimination.}
\label{tab:spt6_mutant}
\begin{tabular}{lccc}
\toprule
\textbf{Genotype} & \textbf{lncRNA Production} & \textbf{DNA Elimination} & \textbf{Progeny Viability} \\
\midrule
Wild-type Spt6 & Normal & Normal & Yes \\
SUMO-defective Spt6 & Normal & Normal & Yes \\
EMA2 KO & Abolished & Partially blocked & No \\
\bottomrule
\end{tabular}
\end{table}

\section{Discussion}

Our findings establish Ema2 as a conjugation-specific SUMO E3 ligase that links the SUMOylation machinery to lncRNA transcription from the parental MAC, a critical step in the scnRNA-mediated genome surveillance pathway in \textit{Tetrahymena}. The pathway can be summarized as follows: (1) scnRNAs are produced from the MIC genome during early conjugation; (2) Ema2-dependent SUMOylation promotes Spt6 and RNAPII chromatin association in the parental MAC; (3) genome-wide lncRNA transcription from the parental MAC provides the scanning substrate; (4) TDSD selectively degrades MDS-matching scnRNAs; (5) retained IES-matching scnRNAs guide heterochromatin formation and DNA elimination in the new MAC.

The observation that SUMOylation-defective Spt6 does not phenocopy EMA2 KO is particularly informative. While Spt6 is a bona fide Ema2 SUMOylation target (confirmed by mass spectrometry and slower-migrating bands on western blot), its SUMOylation is neither necessary nor sufficient for lncRNA transcription. This implies that Ema2 acts through multiple SUMOylation substrates, potentially including other chromatin remodelers, histone modifiers, or transcription factors. The identity of these additional targets represents an important future direction.

The conjugation-specific expression and MAC localization of Ema2 ensure that its lncRNA-promoting activity is temporally restricted to the mid-conjugation window when parental MAC scanning must occur. This temporal control may prevent promiscuous lncRNA transcription during vegetative growth, where it would be unnecessary and potentially deleterious to gene expression.

Our work extends the emerging connection between SUMOylation and small RNA-directed silencing pathways. In \textit{Drosophila}, the SUMO E3 ligase Su(var)2-10 connects piRNA target recognition to transcriptional silencing at transposon loci \citep{ninova2020suvar}. In \textit{Tetrahymena}, Ema2 operates at a distinct step---promoting the lncRNA transcription that enables target recognition---but the convergent involvement of SUMO E3 ligases suggests that SUMOylation is a broadly conserved mechanism for coupling chromatin state changes to small RNA pathways.

\section{Methods}

\subsection{Strains and Culture Conditions}

\textit{Tetrahymena thermophila} wild-type strains B2086, CU427, and CU428 were used. MIC-defective star strains B*VI and B*VII were used for specific crosses. Cells were grown in SPP medium with 2\% proteose peptone at 30$^\circ$C. For conjugation, growing cells ($\sim$5--7 $\times$ 10$^5$/mL) were washed, pre-starved for 12--24 hours, and two mating types were mixed in 10 mM Tris pH 7.5 at 30$^\circ$C.

\subsection{Gene Knockout and Tagging}

EMA2 somatic knockout was generated by disrupting all MAC copies via biolistic transformation. The Ema2 protein was C-terminally HA-tagged at the endogenous locus using construct pEMA2-HA-neo3, selected with 100 $\mu$g/mL paromomycin and 1 $\mu$g/mL CdCl$_2$.

\subsection{DNA Elimination Assay}

FISH was performed at 36 hpm using probes against the Tlr1 transposable element. Cells were fixed, hybridized with labeled probes, and imaged by fluorescence microscopy.

\subsection{Northern Blot for scnRNAs}

Total RNA was extracted at 3, 4.5, and 6 hpm. scnRNAs ($\sim$26--31 nt) were separated on denaturing polyacrylamide gels and probed with labeled oligonucleotides complementary to MDS regions (Mi-9 probe) or IES regions.

\subsection{In Vitro Interaction and SUMOylation Assays}

Recombinant GST-Ema2 and His-Ubc9 were expressed in \textit{E.~coli}, purified, and used for GST pulldown and in vitro SUMOylation assays with purified E1, E2, and Smt3 components.

\subsection{Mass Spectrometry for SUMOylation Targets}

His-tagged Smt3 was expressed in wild-type and EMA2 KO backgrounds. His-Smt3-conjugated proteins were purified by Ni-NTA pulldown from conjugating cells at 6 hpm and analyzed by LC-MS/MS with label-free quantification.

\subsection{RT-PCR for lncRNAs}

RT-PCR was performed at 6 hpm using primers spanning MDS-IES junctions. These primers detect lncRNAs that read through IES boundaries but do not amplify processed mRNAs. Genomic DNA was used as a positive control.

\subsection{Immunofluorescence}

Cells were fixed at indicated time points and stained with anti-HA (1:1,000), anti-H3K27me3 (1:1,000), anti-H3K9me3 (1:1,000), anti-Rpb3 (1:1,000), or J2 anti-dsRNA (1:1,000) antibodies. Secondary antibodies were fluorophore-conjugated. Images were acquired by confocal or widefield fluorescence microscopy.

\section{Conclusion}

We have identified Ema2 as a conjugation-specific SUMO E3 ligase that promotes lncRNA transcription from the parental macronucleus in \textit{Tetrahymena}, establishing a mechanistic link between SUMOylation and small RNA-mediated genome surveillance. Ema2 acts through SUMOylation of Spt6 and likely additional chromatin-associated targets to maintain Spt6 and RNA polymerase II on chromatin during the critical mid-conjugation window. Loss of this activity blocks lncRNA accumulation, prevents TDSD, impairs heterochromatin formation, partially blocks DNA elimination of $\sim$12,000 IESs, and abolishes sexual progeny viability. These findings reveal SUMOylation as a conserved mechanism for regulating lncRNA transcription in the context of small RNA-directed chromatin modification.

\bibliographystyle{plainnat}
\bibliography{references}

\end{document}
